% ============================================================
%  Notation, symbol table, sign conventions
% ============================================================

\paragraph{Indices and sets.}
Indices are lower-case Latin letters; sets are calligraphic upper-case:
\begin{itemize}[leftmargin=*, itemsep=0pt]
\item $\rx \in \rxset$ is the receiver.  Single-receiver modes (PPP,
      PPP-RTK) have $|\rxset| = 1$.  Baseline modes have a fixed
      reference and one or more rovers.  The network producer
      (\texttt{PPP\_RTK\_Network}) processes a peer set of
      $|\rxset| \ge 2$ stations (50--305 in practice) with no
      reference/rover asymmetry.
\item $g \in \sysset = \{\text{GPS}, \text{GAL}, \text{BDS},
      \text{GLO}, \text{QZS}, \dots\}$ is the constellation.
\item $\sat \in \satset{g}$ is a satellite of constellation $g$.
      $\piv{g}$ denotes the pivot satellite of constellation $g$
      under the BSSD parameterisation (\Cref{sec:diff}).
\item $\sigl \in \sigset$ is the signal/observation type, identified by
      its RINEX-3 three-character code (e.g.\ \texttt{C1C}, \texttt{L2W},
      \texttt{L5X}).
\item $k \in \mathbb{N}$ is the discrete epoch index.
\end{itemize}

\paragraph{Vectors and matrices.}
Vectors are bold lower-case ($\bm{x}$, $\rxpos$); matrices are bold
upper-case ($\jacH$, $\covM$).  $\bm{a}^{\Tr}$ is the transpose.
A hat denotes a state estimate and a tilde a one-step prediction:
$\hat{\stateV}_{k|k}$ and $\tilde{\stateV}_{k|k-1}$.

\paragraph{Observables.}
Each receiver/satellite/signal triple supplies two real-valued
observables per epoch, both in metres:
\begin{itemize}[leftmargin=*, itemsep=0pt]
\item $\code^\sat_{\rx,\sigl}$ is the pseudorange (code), in metres.
\item $\phaseo^\sat_{\rx,\sigl}$ is the carrier phase, scaled to metres
      by $\lambda_\sigl$.  In addition to the range, clock and atmosphere
      terms it shares with the code, it carries an ambiguous (non-range)
      part: $\lambda_\sigl N^\sat_{\rx,\sigl}$ plus the fractional phase
      biases $b_{\rx,\sigl} - b^\sat_\sigl$ (\Cref{eq:obs-phase}).
\end{itemize}
Where the distinction matters, $\sigl$ also denotes the band
($L_1, L_2, L_5, \ldots$) and the tracking attribute (e.g.\ \texttt{C}, \texttt{W}, \texttt{X}).
The leading observable-type letter (\texttt{C}/\texttt{L}) selects code
or phase.  The frequency-only quantities $f_\sigl$, $\lambda_\sigl$ and
$N^\sat_{\rx,\sigl}$ depend only on the (band, attribute) part, so a code
signal and its sibling phase share them.  The state keys follow the same
split: ambiguity keys carry the full phase code (\texttt{"G01\_L1C"}),
receiver code-bias keys only the band digit (\texttt{"GPS\_IFB2"}).

\paragraph{Frequencies and wavelengths.}
$f_\sigl$ is the carrier frequency (Hz) and $\lambda_\sigl = \sLight / f_\sigl$
the wavelength (m).  The frequency-squared factor that maps the
L1-equivalent slant ionospheric delay onto signal $\sigl$ is
\begin{equation}
\mu_\sigl \;:=\; \frac{f_1^2}{f_\sigl^2}.
\label{eq:iono-scaling}
\end{equation}
Here $f_1$ is the GPS~L1 frequency ($1575.42$~MHz), the common
reference for all constellations; it is not the band-1 frequency of
the satellite's own system.  The estimated state $I^\sat_{\rx,1}$ is
therefore GPS-L1-equivalent for BDS, GLO and GAL satellites as well
(\texttt{MeasurementModel::buildJacobianRow}: \texttt{f\_ref =
Constants::GPS\_L1\_FREQ}).

\paragraph{Differencing operators.}
Three differencing operators recur throughout:
\begin{equation*}
\bsd{q^{\sat,\piv{g}}_{\rx,\sigl}}
  := q^\sat_{\rx,\sigl} - q^{\piv{g}}_{\rx,\sigl},
\quad
\brd{q^{\sat}_{\rx_1\rx_2,\sigl}}
  := q^\sat_{\rx_1,\sigl} - q^\sat_{\rx_2,\sigl},
\quad
\dd{q^{\sat,\piv{g}}_{\rx_1\rx_2,\sigl}}
  := \bsd{\big(\brd{q^\sat_{\rx_1\rx_2,\sigl}}\big)}.
\end{equation*}
$\bsd{q}$ is the between-satellite single difference of $q$
relative to the pivot $\piv{g}$, $\brd{q}$ the between-receiver
single difference, and $\dd{q}$ the double difference.

\paragraph{Sign conventions.}
The slant ionospheric delay $I^\sat_{\rx,\sigl}$ enters the code
observable with a positive sign and the carrier phase observable with a
negative sign; this sign difference produces code--phase divergence.
All other range-additive terms (geometry, clocks, troposphere, biases)
carry the same sign on both observables.
\Cref{eq:obs-code,eq:obs-phase} give the explicit equations.

\paragraph{Symbol table.}
\begin{center}
\begin{tabular}{@{}llp{0.55\linewidth}@{}}
\toprule
Symbol & Units & Meaning \\
\midrule
$\rxpos_\rx, \satpos$         & m   & ECEF position of receiver/satellite \\
$\geomrange^\sat_\rx$         & m   & geometric range $\|\satpos - \rxpos_\rx\|$ \\
$\los^\sat_\rx$               & none & unit LOS from $\rx$ to $\sat$ in ECEF \\
$\sLight\,\mathrm{d}t_\rx$    & m   & receiver clock offset (scaled by $c$) \\
$\sLight\,\mathrm{d}t^\sat$   & m   & satellite clock offset (CLK product in user modes; estimated state \texttt{"<sat>|SC"} in \texttt{PPP\_RTK\_Network}) \\
$\sLight\,\delta t^{(g)}_\rx$ & m   & receiver inter-system time offset for constellation $g\neq g_{\mathrm{ref}}$ \\
$T^\sat_\rx$                  & m   & total tropospheric path delay \\
$\zhd^\sat_\rx, \zwd^\sat_\rx$ & m  & zenith hydrostatic / wet components \\
$I^\sat_{\rx,1}$              & m   & L1-equivalent slant ionospheric delay \\
$N^\sat_{\rx,\sigl}$          & cyc & integer carrier-phase ambiguity \\
$b_{\rx,\sigl}, b^\sat_\sigl$ & m   & receiver / satellite phase biases \\
$\bar b^\sat_\sigl$           & m   & estimable satellite phase bias (network mode; \Cref{sec:state}) \\
$B_{\rx,\sigl}, B^\sat_\sigl$ & m   & receiver / satellite code biases (OSB) \\
$B^{(g,b)}_{\rx,k}$           & m   & per-FCN-slot GLONASS receiver IFB, band $b$; here $k \in [-7,+6]$ is the signed FDMA channel number (a context-local reuse, not the epoch index) \\
$\varepsilon^\code, \varepsilon^\phaseo$ & m & residual (multipath + thermal) \\
$\bm{v}_\rx$                  & m/s & receiver velocity (optional / IMU state) \\
$\lamwl,\ \lamnl,\ \lamewl$   & m   & wide- / narrow- / extra-wide-lane wavelength $\sLight/(f_a \mp f_b)$ for the lane's band pair \\
$\lamif$                      & m   & effective ionosphere-free wavelength $|c_1\lambda_1 - c_2\lambda_2|$ with the IF coefficients $c_1 - c_2 = 1$; not an integer wavelength (\dcref{sec:slip}{the cycle-slip document}) \\
$\Cbe$                        & none & body-to-ECEF attitude (direction-cosine) matrix \\
\bottomrule
\end{tabular}
\end{center}
